\section{Component Evidence and Controls}
\label{app:heads}

\subsection{Complete head comparisons}
Table~\ref{tab:seven-heads} reports correct counts for every configuration in the head comparison. The R1 wrapper uses diagnosis on 1,780 units, support rules on 72, and the saved original answer on 1,014. R2 uses rules on 55 units and preserves 24 original answers. R3 uses diagnosis on 680 units, rules on ten, and preserves 43 original answers. These are unit-level coverage counts for the head experiment; the full controller's input-level coverage additionally includes reference endpoints and other compatible tasks.

% Generated from existing archived head outcomes; correct counts.
\begin{table}[tb]
\centering
\small
\caption{Archived original $\rightarrow$ head correct counts on R1--R3 for all seven configurations. The head combines candidate re-decision with rule support and preserves original answers outside coverage. These are analyses of exposed development data, using saved method-context decisions and shared backbone-level rule facts; they are not seven new upstream pipeline reruns. R3 is contained in R1; the three sets contain 2,941 distinct units.}
\label{tab:seven-heads}
\begin{tabular*}{\linewidth}{@{\extracolsep{\fill}}lrrr@{}}
\toprule
System & R1 ($n=2{,}866$) & R2 ($n=79$) & R3 ($n=733$) \\
\midrule
M0 Direct & $861 \rightarrow 1,329$ & $7 \rightarrow 29$ & $131 \rightarrow 338$ \\
M2 MedRGAG & $898 \rightarrow 1,323$ & $6 \rightarrow 23$ & $181 \rightarrow 340$ \\
M3 MA-RAG & $766 \rightarrow 1,260$ & $27 \rightarrow 43$ & $63 \rightarrow 238$ \\
M4 MedRAG--L & $798 \rightarrow 1,225$ & $23 \rightarrow 39$ & $223 \rightarrow 386$ \\
M5 MedRAG--Q & $819 \rightarrow 1,367$ & $33 \rightarrow 46$ & $185 \rightarrow 443$ \\
M6 MedGENIE & $858 \rightarrow 1,317$ & $15 \rightarrow 38$ & $157 \rightarrow 305$ \\
M7 TC-RAG & $901 \rightarrow 1,331$ & $23 \rightarrow 37$ & $212 \rightarrow 349$ \\
\bottomrule
\end{tabular*}
\end{table}

% Generated from archived controls; all cells are correct counts.
\begin{table}[tb]
\centering
\small
\caption{Controls expose limits of attributing the head gains. (a) R3 diagnosis subset ($n=680$): original, one catalog order (order 0), majority across catalog orders without calibration, and calibrated majority. Calibration does not improve every system. (b) R2 ($n=79$): the same-coverage None control predicts None on 55 covered inputs and preserves 24 original answers. Its higher aggregate score reflects 46 covered None targets; it misses all nine covered cases with remaining support. The last two columns isolate these nine cases.}
\label{tab:head-controls}
\textbf{(a) R3: correct out of 680}\par\smallskip
\begin{tabular*}{\linewidth}{@{\extracolsep{\fill}}lrrrr@{}}
\toprule
System & Original & Order 0 & Majority & Calibrated \\
\midrule
M0 Direct & 116 & 299 & 314 & 314 \\
M2 MedRGAG & 163 & 320 & 325 & 315 \\
M3 MA-RAG & 33 & 215 & 200 & 204 \\
M4 MedRAG--L & 203 & 339 & 359 & 359 \\
M5 MedRAG--Q & 160 & 382 & 402 & 410 \\
M6 MedGENIE & 137 & 296 & 288 & 276 \\
M7 TC-RAG & 186 & 312 & 324 & 318 \\
\bottomrule
\end{tabular*}
\medskip
\textbf{(b) R2: aggregate and remaining-support subset}\par\smallskip
\begin{tabular*}{\linewidth}{@{\extracolsep{\fill}}lrrrr@{}}
\toprule
System & Head /79 & None /79 & Head /9 & None /9 \\
\midrule
M0 Direct & 29 & 48 & 6 & 0 \\
M2 MedRGAG & 23 & 47 & 4 & 0 \\
M3 MA-RAG & 43 & 58 & 3 & 0 \\
M4 MedRAG--L & 39 & 58 & 4 & 0 \\
M5 MedRAG--Q & 46 & 56 & 6 & 0 \\
M6 MedGENIE & 38 & 52 & 5 & 0 \\
M7 TC-RAG & 37 & 54 & 2 & 0 \\
\bottomrule
\end{tabular*}
\end{table}


\subsection{Diagnosis: catalog, calibration, and source checks}
The public catalog contains 49 diagnoses, while the frozen task scorer recognizes 46 normalized diagnosis keys. The fixed public catalog is supplied to the decision component independently of the current item's gold answer. On 680 diagnostic units, Table~\ref{tab:head-controls}(a) separates the single-catalog improvement from repeated selection and calibration. Across repaired R3 outcomes, 1,186 of 1,272 previously failed outputs were unmapped diagnoses, 23 were otherwise invalid, and 63 were valid wrong answers. This breakdown indicates that output grounding accounts for a large part of the measured improvement. The task scorer treats unmapped outputs as incorrect.

Additional controls explored alternative name/code presentations on 40 diagnostic inputs. Direct code selection with a 384-token budget obtained 18/21 correct for Llama/Qwen; direct name selection at a 128-token budget obtained 15/18. Reasoning with names at 384 tokens obtained 18/19, while reasoning with codes at that budget obtained 21/22. The direct-name/direct-code comparison jointly varies representation and token budget. The evidence ledger retains the complete run identities and budgets.

Two additional source checks assessed catalog compatibility. On 47 ASCENT records, the calibrated head obtained 38/39 correct for Llama/Qwen, compared with 39/39 for uncalibrated voting and 38/39 for a single catalog decision. These records included 42 source-training and five source-test examples and subsequently informed development. Across three CUPCase batches totaling 18 records, the calibrated head obtained 15/16, uncalibrated voting 16/17, single selection 15/18, and character-length-normalized candidate likelihood 15/14. Catalog-based selection performs well on these additional records; calibration provides no consistent incremental advantage.

\subsection{Support removal and remaining support}
The constant ``None'' control in Table~\ref{tab:head-controls}(b) uses the same 55-input rule coverage and retains the same 24 original answers outside coverage. Forty-six covered targets have ``None'' as their answer; the other nine retain support for at least one option. On those nine, original correct counts for Direct, MedRGAG, MA-RAG, MedRAG-Llama, MedRAG-Qwen, MedGENIE, and TC-RAG are 5/4/1/3/0/2/3, and head counts are 6/4/3/4/6/5/2. This subset measures the ability to retain independent support after another basis is removed.

The 123-input add/remove ablation in the main text holds the extracted facts and statuses fixed. It evaluates the shared executor on the controller's complete rule-input panel, which extends beyond the 79-unit R2 category. The full executor leaves one Llama input and 14 Qwen inputs unavailable; these remain included in the 123-input denominator. The contrast tests the effect of the specified operations under the same upstream extraction.

\subsection{Supplied-model consequence execution}
\label{app:model}
The separate model panel uses 24 virtual parameter and initial-state profiles, each with 12 questions. A question supplies the physiological mechanism, factual schedule, requested local intervention, and output projection. The typed-patch procedure edits the named action and executes the resulting trajectory. The full-plan control has the same tool access but regenerates a complete plan; the no-intervention control retains patch structure while omitting the actual intervention. Table~\ref{tab:historical_model} reports binary, numerical, and full-output correctness for the 288 structured queries over these virtual profiles.

\begin{table}[htbp]
\centering
\small
\caption{Historical explicit-model validation. Each configuration contains 288 questions from 24
virtual parameter/initial-state profiles. Counts report correct outputs within the
supplied model contract.}
\label{tab:historical_model}
\setlength{\tabcolsep}{5pt}
\begin{tabularx}{\linewidth}{@{}lXrrr@{}}
\toprule
Config. & Historical arm & \shortstack{Binary\\/288} & \shortstack{Numerical\\/288} & \shortstack{Full output\\/288} \\
\midrule
M4 & Ordinary reanswer & 118 & 0 & 0 \\
M4 & Full execution & 286 & 270 & 270 \\
M4 & Patch without intervention execution & 234 & 67 & 67 \\
M4 & Typed patch and execution & 288 & 288 & 288 \\
\addlinespace
M5 & Ordinary reanswer & 187 & 0 & 0 \\
M5 & Full execution & 265 & 222 & 222 \\
M5 & Patch without intervention execution & 234 & 67 & 67 \\
M5 & Typed patch and execution & 288 & 288 & 288 \\
\bottomrule
\end{tabularx}
\end{table}


\subsection{Stable readout and generated aggregation}
The F component panel comprises an initial 64 groups, an additional 128 groups, and a final frozen extension of 16 groups. Group-weighted paired scores for the complete 208-group panel are reported in Section~\ref{sec:conditional}. The ordinary branch aggregates valid native answer keys; the NLI branch combines direct and analysis-conditioned semantic probabilities across both label mappings. The component's exact stopping, tie-breaking, and invalid-output policies are specified in Section~\ref{sec:method} and the accompanying implementation.

The development study also tested generated aggregation over the available expert outputs on 121 inputs. The adopted controller obtained 87/82 correct for Llama/Qwen, a three-draw alternative 84/74, and the default generated aggregator 75/73. Adding a newly developed support-addition proposal to the candidate pool produced 72/57 under top-one selection and 79/74 under generated aggregation. These experiments led to the shared support kernel and sparse adopted controller. The 121-input panel serves as a development control separate from the full-v5 evaluation.
